\section{Architecture and Technology Stack}

\subsection{Design Methodology}

The system was developed iteratively, one pipeline stage at a time, each stage
operating end to end before the next was begun. Two principles governed most of the
design decisions:

\begin{itemize}
  \item \textbf{Measurement precedes commitment.} Every threshold and sampling choice
        in the pipeline was established by executing it against real footage and
        comparing outcomes, rather than by selecting a plausible value. Several of the
        enhancements listed in Section~4 exist because the obvious choice was adopted
        first and observed to fail.
  \item \textbf{Extension by registration, not by modification.} Analyzers and
        aggregators are declared in a single registry and discovered from it. No part
        of ingestion, storage or the API names a specific analyzer, so introducing one
        does not require any of them to be edited.
\end{itemize}

Verification was performed against a running server rather than by code inspection,
because a number of defects (stale processes, broken answer synthesis, incorrect
media paths) were observable only over HTTP.

\subsection{Architectural Overview}

The system is a three-tier application. A Next.js client provides the operator
interface and hosts an LLM agent that translates conversational requests into retrieval
calls. A FastAPI service owns the entire video pipeline and exposes it as a
self-describing HTTP API. Persistence is divided across object storage for media, a
relational database for application state, and a vector database for retrieval.

Composition within the analysis service is deliberately in-process: aggregators
consume analyzer output directly rather than over HTTP, and HTTP is reserved for the
external boundary. Analyzers and aggregators are held in registries, which is what
makes NFR-8 achievable, because the ingest path, the store and the API iterate over a
registry and never name a specific analyzer.

The analysis service has no concept of users, projects or row-level security: any
caller able to reach it can read every recording it holds. Multi-tenancy is therefore
enforced entirely in the client tier, by row-level security on the database and by a
single scoping function through which every agent tool resolves the recording
identifiers it is permitted to pass downward (NFR-9).

\begin{table}[htbp]
\centering
\small
\begin{tabularx}{\textwidth}{@{}L{3.2cm}X@{}}
\toprule
\textbf{Layer} & \textbf{Responsibility} \\
\midrule
Client            & Marketing site and documentation, authentication, project workspace, video upload and status, agent chat, clip artifact panel, video studio timeline, and the per-recording detail view. \\
Agent layer       & Multi-step streaming tool loop. Converts a natural-language request into search, question, inspection and clip operations, and streams results and artifacts back to the client. \\
API route layer   & The client's own server routes: project scoping, row-level security, job reconciliation, and proxying to the analysis service. \\
Core service      & Ingest orchestration, chunking, analyzers, aggregators, vector search, question answering, clip and poster export, job tracking. \\
Storage boundary  & The single module that knows both a storage URL and a local path; everything above it speaks URLs, everything below it receives paths. \\
Persistence       & Object storage for video bytes, relational store for application state, vector store for chunk embeddings, local cache for decoded media and model weights. \\
\bottomrule
\end{tabularx}
\end{table}

\subsection{Module Inventory}

\begin{table}[htbp]
\centering
\small
\begin{tabularx}{\textwidth}{@{}L{1.6cm}L{3.4cm}X@{}}
\toprule
\textbf{Tier} & \textbf{Module} & \textbf{Responsibility} \\
\midrule
Core & \texttt{chunk.py}, \texttt{chunking/} & The three chunking modes and the weighted fusion of boundary evidence. \\
Core & \texttt{boundaries/} & The four detectors: speaker change, silence, hard cut, semantic drift. \\
Core & \texttt{analyzers/} & Six per-chunk passes, each owning its frame sampling, prompt and output shape. \\
Core & \texttt{aggregators/} & Twelve video-level passes over saved records, ordered by declared dependencies. \\
Core & \texttt{extractors/} & Shared primitives: frame readers and samplers, audio decoding, vision--language calls, detector wrappers. \\
Core & \texttt{vectordb/} & Local embedding, the Qdrant point schema with five named vectors, and record rendering. \\
Core & \texttt{api/} & \texttt{core.py} holds the logic; \texttt{app.py} is a thin HTTP layer; \texttt{ask.py} performs routed question answering; \texttt{jobs.py} tracks background ingestion. \\
Core & \texttt{storage.py}, \texttt{paths.py} & The storage boundary and every filesystem path, each overridable by environment variable. \\
Core & \texttt{export.py}, \texttt{poster.py} & Per-chunk MP4 export by a single decode pass, and one poster frame per recording. \\
\midrule
Client & \texttt{app/projects/} & Project list, workspace, upload dialog, media grid, and the per-recording detail view. \\
Client & \texttt{app/agent/} & Chat surface, message and tool renderers, artifact panel, clip reel, and the video studio timeline. \\
Client & \texttt{app/api/agent/tools/} & The nine agent tools and the project scoping function that binds them. \\
Client & \texttt{app/api/videos/} & Ingestion, status, aggregates, timeline, details and re-index routes. \\
Client & \texttt{lib/core/} & Typed client for the analysis service, job reconciliation, chunking configuration and formatting helpers. \\
Client & \texttt{lib/supabase/} & Browser, server and admin clients, session middleware, and the schema migrations. \\
Client & \texttt{content/docs/} & The MDX documentation site: concepts, workspace guides and the full API reference. \\
\bottomrule
\end{tabularx}
\end{table}

\subsection{Ingestion Lifecycle}

Ingestion is asynchronous. The service accepts a file or a link, returns immediately
with a job identifier, and processes the recording on a background thread through the
stages specified in Section~5. The job reports its stage as fetching, chunking,
analyzing, indexing, aggregating or complete, and the client polls it, reconciling the
result into the recording's row. Because jobs are held in memory, a restart mid-ingest
is detected by the client, which either recovers the recording from its persisted
records or fails the row explicitly so that re-indexing becomes the route forward.

\subsection{Persistence}

State is divided across four stores: a relational database for application state, object
storage for media, a document store on disk for analyzer and aggregator output, and a
vector store for retrieval. The division, the schema of each store and the identity that
joins them are specified in Section~7.

\subsection{Technology Stack}

\begin{table}[htbp]
\centering
\small
\begin{tabularx}{\textwidth}{@{}L{3.2cm}X@{}}
\toprule
\textbf{Concern} & \textbf{Technology} \\
\midrule
Frontend        & Next.js 16, React 19, TypeScript 5, Tailwind CSS v4, shadcn/ui on Radix primitives \\
Agent \& models & Vercel AI SDK v6; OpenAI, Google, Groq, Cerebras and xAI providers behind one model registry \\
Client state    & Zustand, TanStack Query \\
Playback        & video.js, hls.js \\
Documentation   & Fumadocs over MDX, served from the same application \\
Auth \& tenancy & Supabase Auth with Postgres row-level security \\
Backend         & Python 3.13, FastAPI, Uvicorn \\
ML runtime      & PyTorch 2.11 (CUDA 13), Transformers \\
Boundary detection & pyannote.audio (speaker change), Silero VAD (silence), PySceneDetect (hard cuts), OpenCLIP ViT-B/32 (semantic drift) \\
Analysis models & Vision--language model for scene, people, object and text description; faster-whisper for transcription; YOLO11 and EasyOCR as detection gates \\
Aggregation models & GLiNER for zero-shot named-entity recognition, a transformer classifier for sentiment, an LLM for narrative aggregates \\
Embeddings      & BGE-small-en-v1.5 via sentence-transformers, executed locally \\
Vector store    & Qdrant, one point per chunk with five named vectors \\
Media I/O       & OpenCV, PyAV, soundfile, Pillow \\
Application data & Supabase: Postgres, Auth and Storage \\
Deployment      & Vercel (client), CUDA-capable host (core service), Supabase Cloud, Qdrant \\
\bottomrule
\end{tabularx}
\end{table}

\optfigure{video-mind-tech.png}{Technology stack across the frontend, backend, storage and deployment tiers.}{1.0}{fig:tech}

\subsubsection*{Justification of Principal Choices}

\begin{itemize}
  \item \textbf{Python for the analysis service.} Every model the system depends upon
        (diarization, CLIP, Whisper, YOLO, the embedding model) is distributed
        as a Python library first. No other runtime would have permitted this number
        of models to be assembled in the time available.
  \item \textbf{FastAPI.} Asynchronous request handling for long-running ingest jobs,
        and generated API documentation, which is material when an LLM agent is one of
        the clients.
  \item \textbf{Qdrant.} Named vectors are the determining factor. Storing several
        vectors per chunk under one point is a first-class capability there, and it is
        what enhancement E-3 rests upon. It also runs embedded, so development
        requires no separate server.
  \item \textbf{A local embedding model.} Embedding is performed on every chunk of
        every recording. Paying an API per embedding would make re-indexing expensive;
        a small local model makes it free.
  \item \textbf{Next.js and the Vercel AI SDK.} Streaming responses, tool calling and
        multi-provider model selection are provided, so the agent interface reduced to
        defining tools rather than building a chat runtime.
  \item \textbf{Supabase.} Postgres, authentication, row-level security and file
        storage from a single service, which together constitute the whole of the
        non-video state.
\end{itemize}

\subsection{System Requirements}

\begin{table}[htbp]
\centering
\small
\begin{tabularx}{\textwidth}{@{}L{3.6cm}X@{}}
\toprule
\textbf{Component} & \textbf{Requirement} \\
\midrule
GPU              & CUDA-capable NVIDIA GPU, 8\,GB VRAM minimum (developed on an RTX 4060). Required for diarization, CLIP, Whisper, YOLO and the embedding model. \\
Runtime          & Python 3.13 for the analysis service; Node.js 20 or later for the client. \\
Memory / storage & 16\,GB system RAM; disk sized for the media cache, which holds one decoded copy of each ingested recording keyed by content hash. \\
External services & Object storage and Postgres (Supabase project); an LLM provider API key for vision analysis and answer synthesis; a Hugging Face token for the gated diarization model. \\
Network          & Outbound HTTPS for model downloads, object storage and LLM calls. Ingestion executes as a background job, so link latency does not block the client. \\
\bottomrule
\end{tabularx}
\end{table}

\subsection{Known Limitations}

The following limitations are known and accepted for the current phase:

\begin{itemize}
  \item \textbf{Entity linking is confined to a single recording.} The clustering step
        generalises to a camera network unchanged, but the available footage shares no
        people, so it could not be validated.
  \item \textbf{Within-chunk person tracking fragments.} Bounding-box association at
        1\,fps loses a person who moves between sampled frames, so a box identifier is
        a referent within a chunk rather than an identity claim.
  \item \textbf{Payload indexes are inert in embedded mode.} Filtering is correct but
        scans. Running Qdrant as a server resolves this with no code change.
  \item \textbf{Jobs are held in memory.} A restart discards job history; records and
        vectors survive on disk, and the client detects and recovers from the lost
        job.
  \item \textbf{Chapters may collapse to one} on unbroken single-location footage.
        This is arguably correct but of little use in the interface.
\end{itemize}

\subsection{Scope for Future Enhancements}

The design reserves room in the areas expected to be extended next:

\begin{itemize}
  \item \textbf{Cross-recording person linking.} People are already linked within a
        recording; performing the same clustering step over a wider set extends this
        across an entire camera network once suitable footage is available.
  \item \textbf{Live camera input.} Ingestion presently accepts a recording. The same
        pipeline applies to a rolling window from a live feed, since chunking operates
        on segments rather than whole files.
  \item \textbf{Visual re-identification.} Where clothing descriptions are too generic
        to separate two people, the stored bounding boxes provide the crops required
        to compare them visually.
  \item \textbf{Additional analyzers.} Any pass that reads a chunk and returns
        structured output, such as licence plates, pose estimation or crowd density,
        may be added as a single module with no change elsewhere.
  \item \textbf{Agent access over MCP.} The API already describes itself; exposing it
        as a tool server would permit an external assistant to use it directly.
\end{itemize}
